\subsection{Numeric Convergence}
\label{sec:numeric_convergence}

We begin by establishing an integration scheme. For this we need to compare two identical expressions
\begin{equation}
\boxed{\begin{aligned}
&\rho^{2}(\mathcal{E}_{F})\int_{0}^{\infty}
 f^{T}(\mathcal{E}+\hbar\omega)\big[1-f^{T}(\mathcal{E})\big]\,\, d\mathcal{E} \\
\; &\rho^{2}(\mathcal{E}_{F}) \frac{k_BT}{e^{\beta\hbar\omega}-1}\Big[\beta\hbar\omega+\ln\left(1+e^{-\beta\mu}\right)-\ln\left(1+e^{\beta(\hbar\omega-\mu)}\right)\Big]
\end{aligned}}
\label{eq:numeric_convergence_reference}
\end{equation}

This serves as a good testing ground for numerical convergence as we expect to be able to reach near machine accuracy. We are going to be using relative error as our comparison metric throughout the following chapter

\subsubsection{Integration Scheme}
Unsurprisingly, for the same integration interval and step size, Simpson integration is much more reliable. I have also tested a quadrature which resulted in a very unstable result.

\begin{figure}[H]
  \centering
  \includegraphics[width=0.8\textwidth]{stage_2_rel_error_grid.png}
  \caption{Figure 3.?: Point-wise relative error evaluated on the default energy-temperature grid}
  \label{fig:stage_2_rel_error_grid}
\end{figure}

Figure 3.? maps the relative error across the $(T, \hbar\omega)$ plane. A distinct "cone" of higher relative error is observed expanding from the Fermi energy with increasing temperature. This behavior is expected because at large photon energies ($\hbar\omega \gg k_B T$), the transition rates are exponentially suppressed. In these regions, the relative error calculation involves dividing by a near-zero reference value, making it highly sensitive to machine epsilon and finite truncation effects. However, in the physically significant spectral region (high signal), the scheme maintains excellent accuracy.

\subsubsection{Step size convergence}
\begin{figure}[H]
  \centering
  \includegraphics[width=0.8\textwidth]{step_size_convergence.png}
  \caption{step\_size\_convergence.}
  \label{fig:step_size_convergence}
\end{figure}
The step convergence plot demonstrates that refining the grid size $\Delta\mathcal{E}$ leads to a rapid exponential reduction in error. A step size of $\Delta\mathcal{E} \approx 10^{-4}$ eV is sufficient to reach the noise floor, where the error becomes dominated by floating-point precision rather than discretization.

\subsubsection{Interval Length Convergence}

\begin{figure}[H]
  \centering
  \includegraphics[width=0.8\textwidth]{interval_length_convergence.png}
  \caption{interval\_length\_convergence.}
  \label{fig:interval_length_convergence}
\end{figure}

The length convergence analysis indicates that the integral value saturates quickly as the upper integration limit is increased. An upper limit of approximately $5.8$ eV relative to the Fermi level is found to be sufficient; extending the range further yields diminishing returns due to the exponential decay of the thermal occupation factors.

\begin{figure}[H]
  \centering
  \includegraphics[width=0.8\textwidth]{stage_2_rel_error_grid_50x50.png}
  \caption{stage\_2\_rel\_error\_grid\_50x50.}
  \label{fig:stage_2_rel_error_grid_50x50}
\end{figure}

With a numerically stable and demonstrably converged treatment of Eq.~\eqref{eq:numeric_convergence_reference} in hand, we have established robust integration parameters ($d\mathcal{E}=10^{-4}$ eV, range limit $\approx 5.8$ eV), as well as written the integrand in a more numerically stable form. We now turn to assess the validity of the physical approximations which will be the focus of the following sections.

\subsubsection{Log-stable thermal factor}
The factor $f^{T}(\mathcal{E}+\hbar\omega)[1-f^{T}(\mathcal{E})]$ is evaluated in log form to avoid overflow/underflow. Let
\begin{equation}
a=\beta(\mathcal{E}-\mu),\qquad b=a+\beta\hbar\omega.
\label{eq:log_stable_definitions}
\end{equation}
Then
\begin{equation}
f^{T}(\mathcal{E}+\hbar\omega)\,[1-f^{T}(\mathcal{E})]
=\frac{e^{a}}{(1+e^{a})(1+e^{b})},
\label{eq:log_stable_fraction}
\end{equation}
and therefore
\begin{equation}
\log\!\Big(f^{T}(\mathcal{E}+\hbar\omega)\,[1-f^{T}(\mathcal{E})]\Big)
=a-\log(1+e^{a})-\log(1+e^{b}).
\label{eq:log_stable_log}
\end{equation}
In code, $\log(1+e^x)$ is computed as $\mathrm{logaddexp}(0,x)$, and the exponentiation is performed once at the end.
